% =====================================================================
\chapter{Problemas}
\label{cap:problemas}
% =====================================================================

Nueve problemas, uno por módulo, con un total de 60 puntos. La puntuación de cada apartado figura entre paréntesis. Los apartados marcados con \emph{Fuente} están adaptados de exámenes reales de cursos universitarios (MIT y Universidad de Cambridge); el enunciado es de redacción propia y los datos se han cambiado cuando convenía para resolverlos a mano. Las soluciones completas están en el \cref{ap:soluciones}.

% ---------------------------------------------------------------------
\section*{Problema B.1. Lógica, contrarrecíproco e inducción (módulo 0.1)}
\phantomsection\addcontentsline{toc}{section}{Problema B.1. Lógica, contrarrecíproco e inducción}

\begin{ejercicio}{3}{eje:logica}
\begin{enumerate}[label=(\alph*)]
  \item \pt{1} Demuestre con una tabla de verdad que $(p\to q)\wedge(p\to r)$ es lógicamente equivalente a $p\to(q\wedge r)$.
  \item \pt{1} Una función $f\colon\R\to\R$ es continua en $a$ si
  \[
  \forall\varepsilon>0\ \exists\delta>0\ \forall x\in\R:\ \abs{x-a}<\delta\ \Rightarrow\ \abs{f(x)-f(a)}<\varepsilon .
  \]
  Escriba la negación de este enunciado de modo que el símbolo $\neg$ no aparezca delante de ningún cuantificador ni de ningún condicional.
  \item \pt{2} Demuestre por contrarrecíproco que, para todo $n\in\Z$, si $n^{2}-6n+5$ es par, entonces $n$ es impar.
  \item \pt{3} Demuestre por inducción que, para todo $n\ge1$,
  \[
  F_{2n}=F_n\,(F_{n-1}+F_{n+1})\qquad\text{y}\qquad F_{2n+1}=F_n^{2}+F_{n+1}^{2},
  \]
  donde $(F_n)$ es la sucesión de Fibonacci con $F_0=0$ y $F_1=1$. (Indicación: demuestre las dos identidades a la vez.)
\end{enumerate}
\fuente{apartado~(d) adaptado de Cambridge, \emph{Mathematical Tripos} Part~IA, Paper~4, 5 de junio de 2019, pregunta~5E(b) (Numbers and Sets)~\cite{tripos19p4}.}
\end{ejercicio}

% ---------------------------------------------------------------------
\section*{Problema B.2. Inyectividad, biyección e inversa (módulo 0.2)}
\phantomsection\addcontentsline{toc}{section}{Problema B.2. Inyectividad, biyección e inversa}

\begin{ejercicio}{3}{eje:funciones}
\begin{enumerate}[label=(\alph*)]
  \item \pt{2} Sea $f\colon\R\setminus\{1\}\to\R\setminus\{2\}$ dada por $f(x)=\dfrac{2x+3}{x-1}$. Compruebe que $f$ está bien definida (es decir, que $f(x)\neq2$), demuestre que es biyectiva y halle su inversa.
  \item \pt{2} Sean $f\colon A\to B$ y $g\colon B\to C$. Demuestre que si $g\circ f$ es sobreyectiva, entonces $g$ es sobreyectiva. Muestre con un contraejemplo que $f$ no tiene por qué ser sobreyectiva.
  \item \pt{3} Sea $f\colon X\to Y$. Demuestre que las tres afirmaciones siguientes son equivalentes:
  \begin{enumerate}[label=(\roman*)]
    \item $f$ es inyectiva;
    \item $f^{-1}\bigl(f(A)\bigr)=A$ para todo $A\subseteq X$;
    \item $f(A\cap B)=f(A)\cap f(B)$ para todos $A,B\subseteq X$.
  \end{enumerate}
\end{enumerate}
\fuente{apartado~(c) adaptado de Cambridge, \emph{Mathematical Tripos} Part~IA, Paper~4, 5 de junio de 2019, pregunta~7E(a) (Numbers and Sets)~\cite{tripos19p4}.}
\end{ejercicio}

% ---------------------------------------------------------------------
\section*{Problema B.3. Sumatorios, desigualdades y logaritmos (módulo 0.3)}
\phantomsection\addcontentsline{toc}{section}{Problema B.3. Sumatorios, desigualdades y logaritmos}

\begin{ejercicio}{2}{eje:sumas}
\begin{enumerate}[label=(\alph*)]
  \item \pt{2} Demuestre, mediante una suma telescópica, que $\displaystyle\sum_{k=1}^{n}\frac{1}{k(k+1)}=\frac{n}{n+1}$ para todo $n\ge1$. Deduzca el valor de la serie $\sum_{k\ge1}\frac{1}{k^{2}+k}$ y diga si es racional.
  \item \pt{2} Demuestre por inducción que $\displaystyle\sum_{k=1}^{n}k^{2}=\frac{n(n+1)(2n+1)}{6}$ para todo $n\ge1$ y calcule $\displaystyle\sum_{k=11}^{20}k^{2}$.
  \item \pt{1} Demuestre que $x+\dfrac1x\ge2$ para todo $x>0$, con igualdad si y solo si $x=1$.
  \item \pt{1} Un tokenizador tiene un vocabulario de $V=50\,000$ tokens. ¿Cuántos bits, como mínimo, necesita un identificador entero sin signo de longitud fija para numerar todos los tokens? ¿Cuánta información, en bits y en nats, aporta un token elegido de manera uniforme en el vocabulario?
\end{enumerate}
\fuente{apartado~(a) adaptado de Cambridge, \emph{Mathematical Tripos} Part~IA, Paper~4, 5 de junio de 2019, pregunta~2E (Numbers and Sets)~\cite{tripos19p4}.}
\end{ejercicio}

% ---------------------------------------------------------------------
\section*{Problema B.4. Identidades trigonométricas y números complejos (módulo 0.4)}
\phantomsection\addcontentsline{toc}{section}{Problema B.4. Identidades trigonométricas y números complejos}

\begin{ejercicio}{3}{eje:complejos}
\begin{enumerate}[label=(\alph*)]
  \item \pt{2} Demuestre que, para todo $\theta$ que no sea múltiplo entero de $\pi/2$,
  \[
  \frac{\sin3\theta}{\sin\theta}-\frac{\cos3\theta}{\cos\theta}=2 .
  \]
  \item \pt{2} Usando el teorema de De Moivre, demuestre que $\cos3\theta=4\cos^{3}\theta-3\cos\theta$ y que $\sin3\theta=3\sin\theta-4\sin^{3}\theta$ para todo $\theta\in\R$.
  \item \pt{2} Enuncie el teorema de De Moivre y úselo para deducir la fórmula de las raíces $n$-ésimas de un número complejo no nulo $w=\rho e^{i\varphi}$. Calcule en forma binómica las raíces cúbicas de $8i$ y compruebe que su suma es~$0$.
  \item \pt{1} Calcule $\dfrac{(1+i)^{2025}}{(1-i)^{2023}}$.
\end{enumerate}
\fuente{apartado~(c) adaptado de Cambridge, \emph{Mathematical Tripos} Part~IA, Paper~1, 30 de mayo de 2013, pregunta~1C(a)~\cite{tripos13p1}; apartado~(d) adaptado del mismo examen de 30 de mayo de 2019, pregunta~1C(b)~\cite{tripos19p1} (Vectors and Matrices).}
\end{ejercicio}

% ---------------------------------------------------------------------
\section*{Problema B.5. Límites, derivada y teorema del valor medio (módulo 0.5)}
\phantomsection\addcontentsline{toc}{section}{Problema B.5. Límites, derivada y teorema del valor medio}

\begin{ejercicio}{3}{eje:tvm}
\begin{enumerate}[label=(\alph*)]
  \item \pt{2} Demuestre con la definición $\varepsilon$--$\delta$ que $\displaystyle\lim_{x\to2}(x^{2}+x)=6$.
  \item \pt{1} Calcule, a partir de la definición de derivada como límite del cociente incremental, la derivada de $f(x)=\dfrac{x}{x+1}$ en todo $x\neq-1$.
  \item \pt{3} Demuestre, usando el teorema del valor medio, que
  \[
  \frac{x}{1+x^{2}}<\arctan x<x\qquad\text{para todo }x>0 .
  \]
  \item \pt{1} Calcule $\displaystyle\lim_{x\to0}\frac{x-\arctan x}{x^{3}}$.
\end{enumerate}
\fuente{apartado~(b) adaptado de MIT 18.01SC, \emph{Single Variable Calculus}, otoño de 2010, examen final, problema~3(a); apartado~(c) adaptado del mismo examen, problema~15 (pregunta adicional)~\cite{mit1801final}.}
\end{ejercicio}

% ---------------------------------------------------------------------
\section*{Problema B.6. Integración y series de potencias (módulo 0.6)}
\phantomsection\addcontentsline{toc}{section}{Problema B.6. Integración y series de potencias}

\begin{ejercicio}{2}{eje:integrales}
\begin{enumerate}[label=(\alph*)]
  \item \pt{1{,}5} Calcule $\displaystyle\int_{1}^{e}x\ln x\,\dd x$.
  \item \pt{1{,}5} Estudie la convergencia de $\displaystyle\int_{1}^{\infty}\frac{\ln x}{x^{2}}\,\dd x$ y, si converge, calcule su valor.
  \item \pt{2} Obtenga la serie de Taylor de $\ln(1+x)$ centrada en~$0$ y su radio de convergencia. Use los dos primeros términos no nulos para aproximar $\ln 1{,}2$.
  \item \pt{1} Acote el error de la aproximación de~(c) con el resto de Lagrange y compruebe la cota con el valor $\ln1{,}2=0{,}18232$.
\end{enumerate}
\fuente{apartado~(a) adaptado de MIT 18.01SC, \emph{Single Variable Calculus}, otoño de 2010, examen final, problema~8(b); apartados~(c) y~(d) adaptados del mismo examen, problema~14, que usa $\ln(3/2)$~\cite{mit1801final}.}
\end{ejercicio}

% ---------------------------------------------------------------------
\section*{Problema B.7. Gradiente, extremos y multiplicadores de Lagrange (módulo 0.7)}
\phantomsection\addcontentsline{toc}{section}{Problema B.7. Gradiente, extremos y multiplicadores de Lagrange}

\begin{ejercicio}{2}{eje:varias}
\begin{enumerate}[label=(\alph*)]
  \item \pt{1} Sea $f(x,y)=x^{2}y^{3}-2x$. Calcule $\nabla f(2,1)$ y la ecuación del plano tangente a la gráfica de $f$ en el punto $(2,1,f(2,1))$.
  \item \pt{1} Use la aproximación lineal para estimar $f(2{,}1;\,0{,}9)$ y compárela con el valor exacto. Calcule la derivada direccional de $f$ en $(2,1)$ en la dirección del vector $(3,4)$ y el valor máximo de la derivada direccional en ese punto.
  \item \pt{2} Halle y clasifique los puntos críticos de $g(x,y)=x^{3}-3x+y^{2}-4y$.
  \item \pt{2} Halle, con multiplicadores de Lagrange, los valores máximo y mínimo de $h(x,y)=xy$ sobre la elipse $x^{2}+4y^{2}=8$, y los puntos en que se alcanzan. Justifique que existen.
  \item \pt{1} Calcule $\displaystyle\iint_{D}e^{-(x^{2}+y^{2})}\,\dd A$, donde $D$ es el disco $x^{2}+y^{2}\le4$.
\end{enumerate}
\fuente{apartados~(a) y~(b) adaptados de MIT 18.02SC, \emph{Multivariable Calculus}, otoño de 2010, examen~2, problema~1; apartado~(c) adaptado del mismo examen, problema~3(a)~\cite{mit1802ex2}.}
\end{ejercicio}

% ---------------------------------------------------------------------
\section*{Problema B.8. Conteo, Bayes y esperanza (módulo 0.8)}
\phantomsection\addcontentsline{toc}{section}{Problema B.8. Conteo, Bayes y esperanza}

\begin{ejercicio}{2}{eje:probabilidad}
\begin{enumerate}[label=(\alph*)]
  \item \pt{1} ¿De cuántas maneras pueden elegirse 4 tokens distintos de un vocabulario de 12 si importa el orden? ¿Y si no importa?
  \item \pt{1} Demuestre mediante doble conteo la identidad de Pascal $\binom{n}{k}=\binom{n-1}{k-1}+\binom{n-1}{k}$ para $1\le k\le n-1$.
  \item \pt{2} Una bolsa contiene tres monedas cuyas probabilidades de cara son $0{,}5$, $0{,}7$ y $0{,}2$. Se elige una moneda al azar y se lanza tres veces. Calcule la probabilidad de obtener, en este orden, cara, cruz y cruz. Sabiendo que ha salido esa secuencia, ¿cuál es la probabilidad de que la moneda elegida sea la equilibrada?
  \item \pt{2} En el experimento de~(c), sea $X$ el número de caras en los tres lanzamientos. Calcule $\mathbb E[X]$ y $\PP(X=3)$.
\end{enumerate}
\fuente{apartados~(c) y~(d) adaptados de MIT 18.05, \emph{Introduction to Probability and Statistics}, primavera de 2014, examen~1, problema~3, con monedas de probabilidades $0{,}5$, $0{,}6$ y $0{,}1$~\cite{mit1805ex1}.}
\end{ejercicio}

% ---------------------------------------------------------------------
\section*{Problema B.9. Verosimilitud, intervalos y contrastes (módulo 0.9)}
\phantomsection\addcontentsline{toc}{section}{Problema B.9. Verosimilitud, intervalos y contrastes}

\begin{ejercicio}{2}{eje:estadistica}
\begin{enumerate}[label=(\alph*)]
  \item \pt{2} Una variable de Pareto con parámetro $\alpha>0$ tiene densidad $f(x;\alpha)=\alpha\,x^{-\alpha-1}$ para $x\ge1$ (y $0$ en otro caso). Con la muestra $2$, $4$, $8$, halle el estimador de máxima verosimilitud de $\alpha$ y compruebe que es un máximo.
  \item \pt{2} Se toman $n=16$ datos de una población $\Nor(\mu,4^{2})$, con $\sigma=4$ conocida, y se obtiene $\xbar=3{,}9$. Se contrasta $H_0\colon\mu=2$ frente a $H_1\colon\mu\neq2$ al nivel $\alpha=0{,}05$. Dé la región de rechazo en función de $\xbar$, el estadístico $z$, el $p$-valor y la decisión.
  \item \pt{2} Las latencias, en milisegundos, de 8 consultas a un servidor son $12$, $15$, $11$, $14$, $13$, $16$, $12$, $11$. Suponiendo normalidad, calcule la media muestral, la desviación típica muestral y un intervalo de confianza del 95\,\% para la latencia media.
  \item \pt{1} Se comparan dos modelos en 20 pruebas independientes, cada una con un contraste al nivel $0{,}05$. Si en todas es cierta la hipótesis nula, ¿cuál es la probabilidad de obtener al menos un falso positivo? ¿Qué nivel debe usarse en cada prueba para que esa probabilidad sea como máximo $0{,}05$ según Bonferroni?
\end{enumerate}
\fuente{apartado~(a) adaptado de MIT 18.05, \emph{Introduction to Probability and Statistics}, primavera de 2014, examen~2, problema~2, con la muestra $5$, $2$, $3$; apartado~(b) adaptado del mismo examen, problema~6, con $\xbar=1{,}5$~\cite{mit1805ex2}.}
\end{ejercicio}
